% Preamble
\documentclass[11pt, aspectratio=169]{beamer}
	% Packages
	\usepackage[english]{babel}
	\usepackage{lipsum}
	\usepackage{mathptmx}
	\usepackage{xcolor}
	% Themes
	\usefonttheme[onlymath]{serif}
	\usetheme[nofirafonts]{focus}
	\renewcommand{\thempfootnote}{\arabic{mpfootnote}}
	% Cover
	\title{Microeconometrics Using Stata}
	\subtitle{Data management and graphics: Exercises}
	\author{
		Jhon R. Ordoñez
		\footnote{\label{1} National University of San Cristóbal de Huamanga}
		\footnote{\label{2} Faculty of Economic, Administrative and Accounting Sciences}
		\footnote{\label{3} Professional School of Economics}
	}
	\date{\today}
% Body
\begin{document}
	% Cover
	\begin{frame}
		\maketitle
	\end{frame}
	% Outline
	\begin{frame}[t]
		\frametitle{Outline}
		\tableofcontents
	\end{frame}
	% Section 1
	\section{Exercise 1}
		% Slide 1
		\begin{frame}
			\frametitle{Exercise 1}
				\begin{block}{Exercise 1}
					Type the command \textcolor{blue}{\textbf{display \%10.5f 123.321}}. Compare the results with
					those you obtain when you change the format \textcolor{blue}{\textbf{\%10.5f}} to, respectively,
					\textcolor{blue}{\textbf{\%10.5e}}, \textcolor{blue}{\textbf{\%10.5g}}, \textcolor{blue}{\textbf{\%-10.5f}}, 
					and \textcolor{blue}{\textbf{\%10,5f}} and when you do not specify a format.
				\end{block}
		\end{frame}
	% Section 2
	\section{Exercise 2}
		% Slide 1
		\begin{frame}
			\frametitle{Exercise 2}
			\begin{block}{Exercise 2}
				Consider the example of \textcolor{red!50}{\textbf{section 2.3}} except with the variables
				reordered. Specifically, the variables are in the order \textcolor{blue}{\textbf{age}}, \textcolor{blue}{\textbf{name}},
				\textcolor{blue}{\textbf{income}}, and \textcolor{blue}{\textbf{female}}.
				The three observations are \textcolor{blue}{\textbf{29  “Barry" 40.990 0}}; \textcolor{blue}{\textbf{30  “Carrie" 37.000 1}}; and
				\textcolor{blue}{\textbf{31  “Gary" 48.000 0}}. Use \textcolor{blue}{\textbf{input}} to
				read these data, along with names, into Stata, and list the results. Use a
				text editor to create a comma-separated values file that includes
				variable names in the first line, read this file into Stata by using \textcolor{blue}{\textbf{import}}
				\textcolor{blue}{\textbf{delimited}}, and list the results. Then, drop the first line in the text file,
				read in the data by using \textcolor{blue}{\textbf{import delimited}} with variable names
				assigned, and list the results. Finally, replace the commas in the text
				file with blanks, read the data in by using \textcolor{blue}{\textbf{infix}}, and list the results.
			\end{block}
		\end{frame}
	% Section 3
	\section{Exercise 3}
		% Slide 1
		\begin{frame}
			\frametitle{Exercise 3}
			\begin{block}{Exercise 3}
				Consider the dataset in \textcolor{red!50}{\textbf{section 2.4.}} The \textcolor{blue}{\textbf{er32049}} 
				variable is the last known marital status. Rename this variable as \textcolor{blue}{\textbf{marstatus}}, give the
				variable the label “marital status”, and tabulate \textcolor{blue}{\textbf{marstatus}}. From the
				codebook, marital status is married (1), never married (2), widowed
				(3), divorced or annulment (4), separated (5), not answered or do not
				know (8), and no marital history collected (9). Set \textcolor{blue}{\textbf{marstatus}} to
				missing where appropriate. Use \textcolor{blue}{\textbf{label define}} and \textcolor{blue}{\textbf{label values}} to
				provide descriptions for the remaining categories, and tabulate
				\textcolor{blue}{\textbf{marstatus}}. Create a binary indicator variable equal to 1 if the last
				known marital status is married and equal to 0 otherwise, with
				appropriate handling of any missing data. Provide a summary of
				earnings by marital status. Create a set of indicator variables for
				marital status based on \textcolor{blue}{\textbf{marstatus}}. Create a set of variables that
				interact these marital status indicators with earnings.
			\end{block}
		\end{frame}
	% Section 4
	\section{Exercise 4}
		% Slide 1
		\begin{frame}
			\frametitle{Exercise 4}
			\begin{block}{Exercise 4}
				Consider the dataset in \textcolor{red!50}{\textbf{section 2.6.}} Create a box-and-whisker plot of
				\textcolor{blue}{\textbf{earnings}} (in levels) for all the data and for each year of educational
				attainment (use variable \textcolor{blue}{\textbf{education}}). Create a histogram of \textcolor{blue}{\textbf{earnings}}
				(in levels) using $100$ bins and a kernel density estimate. Do earnings in
				levels appear to be right skewed? Create a scatterplot of \textcolor{blue}{\textbf{earnings}}
				against \textcolor{blue}{\textbf{education}}. Provide a single figure that uses \textcolor{blue}{\textbf{scatterplot}},
				\textcolor{blue}{\textbf{lfit}}, and \textcolor{blue}{\textbf{lowess}} of \textcolor{blue}{\textbf{earnings}} against \textcolor{blue}{\textbf{education}}. Add titles for the
				axes and graph heading.
			\end{block}
		\end{frame}
	% Section 5
	\section{Exercise 5}
		% Slide 1
		\begin{frame}
			\frametitle{Exercise 5}
			\begin{block}{Exercise 5}
				 Consider the dataset in \textcolor{red}{\textbf{section 2.6.}} Create kernel density plots for
				lnearns using the kernel(epan2) option with kernel 
					\begin{equation*}
						K(z) = \dfrac{3}{5}\left(1 - \dfrac{z^{2}}{5}\right)
					\end{equation*}
				for $|z| < 1$ and using the \textcolor{blue}{\textbf{kernel(rectangle)}} option with kernel
				$K(z) = 1/2$ for $|z| < 1$. Repeat with the bandwidth increased from the default to 0.3. What
				makes a bigger difference, choice of kernel or choice of bandwidth?
				The comparison is easier if the four graphs are saved using the
				\textcolor{blue}{\textbf{saving()}} option and then combined using the \textcolor{blue}{\textbf{graph combine}} command.
			\end{block}
		\end{frame}
	% Section 6
	\section{Exercise 6}
		% Slide 1
		\begin{frame}
			\frametitle{Exercise 6}
			\begin{block}{Exercise 6}
				Consider the dataset in \textcolor{red!50}{\textbf{section 2.6.}} For each of the available kernels
				that can be used with the \textcolor{blue}{\textbf{kdensity}} command, obtain a kernel density
				plot for \textcolor{blue}{\textbf{lnearns}} using the default bandwidth, and save the graph using
				the \textcolor{blue}{\textbf{saving()}} option. Then, combine all graphs on one page using the
				\textcolor{blue}{\textbf{graph combine}} command and options such as \textcolor{blue}{\textbf{rows(4) ysize(8) xsize(5)}}.
				Comment on the relative smoothness of the various graphs.
			\end{block}
		\end{frame}
	% Section 7
	\section{Exercise 7}
		% Slide 1
		\begin{frame}
			\frametitle{Exercise 7}
			\begin{block}{Exercise 7}
				Consider the dataset in \textcolor{blue}{\textbf{section 2.6.}} Perform lowess regression of
				\textcolor{blue}{\textbf{lnearns}} on \textcolor{blue}{\textbf{hours}} using the default bandwidth and using bandwidth of
				$0.01$. Does the bandwidth make a difference? A moving average $y$ of 
				after data are sorted by $x$ is a simple case of nonparametric regression
				of $y$ on $x$ . Sort the data by \textcolor{blue}{\textbf{hours}}. Create a centered $25$-period moving
				average of \textcolor{blue}{\textbf{lnearns}} with $i$th observation
					\begin{equation*}
						yma_{i} = \dfrac{1}{25} \sum\limits_{j=-12}^{j=12}y_{i+j}
					\end{equation*}
				This is easiest using \textcolor{blue}{\textbf{forvalues}}. Plot this moving average against
				\textcolor{blue}{\textbf{hours}} using the \textcolor{blue}{\textbf{twoway connected}} graph command. Compare with the
				lowess plot.
			\end{block}
		\end{frame}
	% References
	\begin{frame}[t,allowframebreaks]
		\frametitle{References}
		\bibliography{biblio.bib}
		% Style
		\bibliographystyle{apalike}
		% Authors
		\nocite{cameron2022-I}
		\nocite{cameron2022-II}
	\end{frame}
	% Cover
	\begin{frame}
		\maketitle
	\end{frame}
\end{document}